% !TeX program = xelatex
\documentclass[twoside, 10pt]{ctexart}
\usepackage{geometry}
\geometry{a4paper, left=2.5cm, right=2.5cm, top=2.8cm, bottom=2.5cm}
\usepackage{amsmath, amssymb, amsthm}
\usepackage{graphicx}
\usepackage{booktabs}
\usepackage{algorithm}
\usepackage{algorithmic}
\usepackage{subcaption}
\usepackage{cite}
\usepackage{hyperref}
\hypersetup{colorlinks=true, linkcolor=blue, citecolor=blue, urlcolor=blue}
\usepackage{fancyhdr}
\usepackage{ragged2e} % 用于支持 \justify 命令
\usepackage{lipsum} % 仅用于示例，实际可删除
\usepackage{titlesec} % 用于修改章节标题格式

% 字体设置
\setCJKmainfont[BoldFont=SimHei, ItalicFont=KaiTi]{SimSun} % 正文：五号宋体
\setCJKsansfont{SimHei} % 黑体
\setCJKmonofont{SimSun} % 等宽字体
\setCJKfamilyfont{kaishu}{KaiTi_GB2312} % 楷体_GB2312
\newcommand{\kaishu}{\CJKfamily{kaishu}} % 定义楷体命令
\setmainfont{Times New Roman} % 英文和数字使用Times New Roman

% 字号设置
\newcommand{\xiaowu}{\fontsize{9pt}{11pt}\selectfont} % 小五
\newcommand{\sanhao}{\fontsize{16pt}{18pt}\selectfont} % 三号
\newcommand{\wuhhao}{\fontsize{10.5pt}{12pt}\selectfont} % 五号
\newcommand{\xiaosi}{\fontsize{12pt}{14pt}\selectfont} % 小四

% 设置正文字体为精确的10.5pt
\AtBeginDocument{\fontsize{10.5pt}{12pt}\selectfont}

% 设置章节标题格式：黑体小四，左对齐
\titleformat{\section}[hang]
  {\xiaosi \bfseries}
  {\thesection}
  {1em}
  {}

\pagestyle{fancy}
\fancyhf{}
\fancyhead[RE,LO]{\small 通信学报}
\fancyhead[LE,RO]{\small \thepage}
\renewcommand{\headrulewidth}{0.4pt}

\newcommand{\upcite}[1]{\textsuperscript{[#1]}}

\newtheorem{theorem}{定理}[section]
\newtheorem{lemma}[theorem]{引理}
\newtheorem{definition}{定义}[section]

\renewcommand{\algorithmicrequire}{\textbf{输入:}}
\renewcommand{\algorithmicensure}{\textbf{输出:}}

\begin{document}

\begin{center}
    % 中文标题 - 三号黑体
    {\sanhao \textbf{多用户MIMO系统中基于深度学习的预编码与功率分配方法}}\par
    \vspace{1em}
    % 作者 - 五号楷体_GB2312
    {\wuhhao \kaishu 秦政$^{1}$}\par
    \vspace{0.3em}
    {\small （1. 重庆邮电大学，重庆 400065）}\par
    \vspace{1em}
\end{center}

    % 中文摘要 - 小五字体
    \noindent {\xiaowu \textbf{摘\quad 要}：大规模多输入多输出（MIMO）技术是5G/6G通信系统的关键组成部分，其预编码和功率分配的联合优化是一个非凸、高复杂度的问题。传统迭代算法在用户数增多时收敛缓慢，难以满足低延迟要求。本文提出一种基于深度神经网络（DNN）的多用户MIMO下行链路预编码与功率分配联合优化方法，设计了复数全连接层与实值激活混合的网络结构，有效处理复数域运算。仿真结果表明，所提方法在多种信道场景下均可逼近加权最小均方误差（WMMSE）算法的性能，而计算时间显著降低，且对信道估计误差具有更强的鲁棒性。}\par
    \noindent {\xiaowu \textbf{关键词}：大规模MIMO；深度学习；预编码；功率分配；无监督学习}\par
    \noindent {\xiaowu \textbf{中图分类号}：TN929.5}\par
    \noindent {\xiaowu \textbf{文献标志码}：A}\par
    \noindent {\xiaowu \textbf{DOI}：10.11959/j.issn.1000-436x.20240001}\par
    \vspace{2em}

    % 英文标题
    \begin{center}
    {\Large \textbf{Deep Learning Based Precoding and Power Allocation for Multi-User MIMO Systems}}\par
    \vspace{1em}
    {\large QIN Zheng$^{1}$}\par
    \vspace{0.3em}
    {\small （1. Chongqing University of Posts and Telecommunications, Chongqing 400065, China）}\par
    \vspace{1em}
    \end{center}

    % 英文摘要 - 小五字体
    \noindent {\xiaowu \textbf{Abstract}：Massive multiple-input multiple-output (MIMO) is a cornerstone technology for 5G and beyond communication systems, enabling significant improvements in spectral efficiency and system capacity. However, the joint optimization of precoding and power allocation poses formidable challenges due to its inherent non-convexity and high computational complexity. Traditional iterative approaches, such as the weighted minimum mean square error (WMMSE) algorithm, suffer from slow convergence when scaled to large user populations, hindering real-time implementation. To address these limitations, this paper presents an end-to-end deep learning framework for downlink multi-user MIMO systems. A hybrid neural network architecture is proposed, integrating complex-valued fully connected layers with real-valued activation functions to directly map channel state information to optimal precoding matrices. Comprehensive simulations demonstrate three key contributions: 1) The proposed method achieves within 95% of WMMSE performance across diverse channel conditions; 2) Computational latency is reduced by over two orders of magnitude compared to traditional algorithms; 3) Superior robustness to channel estimation errors is achieved. This work provides a novel data-driven approach for real-time resource optimization in massive MIMO systems.}\par
    \noindent {\xiaowu \textbf{Keywords}：massive MIMO; deep learning; precoding; power allocation; unsupervised learning}\par
    \vspace{2em}

    % 收稿日期上方的横线
    \noindent \rule{3.2cm}{0.4pt}\par
    % 收稿日期、基金项目 - 小五字体，英文数字用Times New Roman
    \noindent {\xiaowu \textbf{收稿日期：}\fontfamily{ptm}\selectfont 2024−03−15}；{\xiaowu \textbf{修回日期：}\fontfamily{ptm}\selectfont 2024−06−20}\par
    \noindent {\xiaowu \textbf{基金项目：}国家自然科学基金资助项目（\fontfamily{ptm}\selectfont No.62271001, No.62001002）}\par
    \noindent {\xiaowu \textbf{Foundation Items:} The National Natural Science Foundation of China (No.62271001, No.62001002)}\par
    \vspace{2em}

\twocolumn

% 允许公式跨页
\allowdisplaybreaks

% 调整表格与上下文的间距
\setlength{\textfloatsep}{1em} % 表格与上下文的间距

\section{引言}

随着移动数据流量的指数级增长，第五代（5G）及未来第六代（6G）移动通信系统对频谱效率和系统容量提出了更高要求。大规模多输入多输出（Massive MIMO）技术通过在基站部署数十至数百根天线，可实现空间复用增益、阵列增益和干扰抑制，已成为提升系统性能的核心技术\upcite{1}。然而，Massive MIMO系统的性能高度依赖于有效的预编码与功率分配策略，其优化问题具有非凸性和高维度特性，给算法设计带来了严峻挑战。

传统预编码方案可分为线性和非线性两类：线性方案如迫零（ZF）、正则化迫零（RZF）等实现简单但性能受限；非线性方案如脏纸编码（DPC）理论性能最优但复杂度极高。近年来，加权最小均方误差（WMMSE）算法通过迭代优化框架在性能与复杂度之间取得了较好平衡，成为业界主流方法\upcite{2}。然而，WMMSE算法的迭代特性导致其计算复杂度随用户数量呈超线性增长，难以满足毫秒级延迟的实时通信需求。

深度学习的兴起为解决复杂优化问题提供了新的思路。文献\upcite{3}率先将全连接网络应用于MIMO预编码设计，但该方法存在输出维度爆炸问题；文献\upcite{4}利用图神经网络建模用户间干扰，实现了可扩展的功率分配。然而，现有深度学习方法普遍存在以下局限：1）将复数信道矩阵拆分为实部和虚部分别处理，破坏了复数运算的整体性；2）过度依赖有监督学习，需要大量高质量标签数据；3）缺乏对信道估计误差的鲁棒性分析。

针对上述问题，本文提出一种基于深度神经网络的端到端预编码与功率分配联合优化方法，主要创新点如下：
\begin{enumerate}
    \item \textbf{复数域神经网络设计}：提出复数全连接层与实值激活函数的混合架构，实现端到端的复数运算，避免实虚部分离带来的性能损失。
    \item \textbf{无监督学习框架}：以负和速率为损失函数，无需预设最优标签，直接优化系统性能指标。
    \item \textbf{鲁棒性分析}：通过理论推导和仿真验证，证明所提方法对信道估计误差具有良好的鲁棒性。
\end{enumerate}

本文余下部分结构如下：第2节建立系统模型并描述优化问题；第3节详细阐述基于深度神经网络的联合优化算法；第4节提供仿真结果与性能分析；第5节总结全文并展望未来研究方向。

\section{系统模型与问题建模}

\subsection{系统模型}

考虑一个多用户MIMO下行链路系统，基站配置$N_t$根发射天线，同时服务$K$个单天线用户。定义$\mathbf{h}_k \in \mathbb{C}^{1 \times N_t}$为用户$k$的信道矢量，$\mathbf{H} = [\mathbf{h}_1^T, \mathbf{h}_2^T, \ldots, \mathbf{h}_K^T]^T \in \mathbb{C}^{K \times N_t}$为聚合信道矩阵。基站发送信号可表示为：

\begin{equation}
\mathbf{x} = \sum_{k=1}^K \mathbf{w}_k s_k
\end{equation}

其中，$\mathbf{w}_k \in \mathbb{C}^{N_t \times 1}$为用户$k$的预编码矢量，$s_k$为归一化数据符号（$\mathbb{E}[|s_k|^2] = 1$）。用户$k$的接收信号为：

\begin{equation}
y_k = \mathbf{h}_k \mathbf{w}_k s_k + \sum_{j \neq k} \mathbf{h}_k \mathbf{w}_j s_j + n_k
\end{equation}

其中，$n_k \sim \mathcal{CN}(0, \sigma^2)$为加性高斯白噪声（假设所有用户噪声功率相同）。定义用户$k$的信干噪比（SINR）为：

\begin{equation}
\gamma_k = \frac{|\mathbf{h}_k \mathbf{w}_k|^2}{\sum_{j \neq k} |\mathbf{h}_k \mathbf{w}_j|^2 + \sigma^2}
\end{equation}

\subsection{优化问题}

本文的优化目标是在满足总功率约束的前提下最大化系统和速率，即：

\begin{subequations}
\begin{align}
\max_{\{\mathbf{w}_k\}_{k=1}^K} \quad & \sum_{k=1}^K \log_2(1 + \gamma_k) \\
\text{s.t.} \quad & \sum_{k=1}^K \|\mathbf{w}_k\|_F^2 \le P_{\text{max}}
\end{align}
\label{eq:optimization}
\end{subequations}

其中，$\|\cdot\|_F$表示Frobenius范数，$P_{\text{max}}$为基站最大发射功率。该问题具有以下特点：

\begin{itemize}
    \item \textbf{非凸性}：目标函数中的对数项和SINR表达式均为非凸函数；
    \item \textbf{耦合性}：各用户的预编码矢量相互耦合，无法独立优化；
    \item \textbf{高维度}：优化变量维度为$K \times N_t$，随用户数和天线数增长而急剧增加。
\end{itemize}

传统WMMSE算法通过交替优化将原问题转化为一系列凸子问题，但迭代收敛速度慢，难以满足实时性要求。本文提出一种基于深度学习的端到端优化框架，直接学习从信道矩阵到预编码矩阵的映射函数$\mathbf{W} = f(\mathbf{H}; \theta)$，其中$\theta$为网络参数。

\section{深度学习预编码与功率分配网络设计}
\subsection{复数神经网络原理}
由于信道矩阵和预编码矢量均为复数，直接使用实数神经网络需将实部虚部分别处理，这会破坏复数相位关系。本文采用复数全连接层\upcite{5}，其权重$\mathbf{W}_c \in \mathbb{C}^{m\times n}$，偏置$\mathbf{b}_c \in \mathbb{C}^{m}$，对复数输入$\mathbf{z} \in \mathbb{C}^{n}$的输出为$\mathbf{y} = \mathbf{W}_c\mathbf{z} + \mathbf{b}_c$。反向传播时，将复数变量视为二维实数向量进行求导。为避免复数激活函数（如modReLU）的不稳定性，隐层后采用CReLU：$\text{CReLU}(z) = \text{ReLU}(\Re(z)) + \mathrm{j} \cdot \text{ReLU}(\Im(z))$。

\subsection{网络结构}
本文设计的深度预编码网络（DPNet）结构如图1所示。输入层接收用户信道矩阵$\mathbf{H}$（展平为$K \times N_t$复数向量）。随后经过三个复数全连接层，神经元数目分别为$512$、$1024$和$2K N_t$。最后一层的输出重新整形为复数预编码矩阵$\mathbf{W} \in \mathbb{C}^{N_t\times K}$。功率约束通过一个投影层实现：
\begin{equation}
\mathbf{W}_{\text{proj}} = \sqrt{P_{\text{max}}} \frac{\mathbf{W}}{\|\mathbf{W}\|_F}
\label{eq:proj}
\end{equation}

\begin{figure}[htbp]
\centering
\includegraphics[width=0.8\linewidth]{example-image}
\caption{DPNet 结构图}
\label{fig:dpnet}
\end{figure}

\subsection{无监督学习策略}
为了直接最大化系统和速率，我们定义损失函数为负和速率：
\begin{equation}
\mathcal{L}(\mathbf{H}) = -\sum_{k=1}^K \log_2\left(1 + \frac{|\mathbf{h}_k \mathbf{w}_k|^2}{\sum_{j\neq k} |\mathbf{h}_k \mathbf{w}_j|^2 + \sigma_k^2}\right)
\end{equation}
注意，$\mathbf{w}_k$是网络输出的预编码矢量。该损失函数关于网络权重可微（通过复数求导链式法则）。训练过程中不需要监督标签（如WMMSE的最优解），仅需随机生成信道样本和噪声功率即可。我们采用Adam优化器，学习率初始为0.001，每50个epoch衰减为原来的0.9倍。

\subsection{与功率分配的关系}
所提出的DPNet直接输出预编码矩阵，其中每个用户的功率分配隐含在$\|\mathbf{w}_k\|^2$中。因此该网络同时实现了预编码方向和功率分配，无需额外模块。

\section{仿真结果与分析}
\subsection{仿真设置}
仿真参数如表1所示。我们将DPNet与以下基准方法对比：1) 迫零预编码 + 等功率分配（ZF-EPA）；2) 迫零预编码 + 注水功率分配（ZF-WF）；3) WMMSE迭代算法（30次迭代作为近似最优）。DPNet使用10万组随机信道进行训练，1万组信道测试。信道模型为瑞利平坦衰落，噪声功率$\sigma_k^2 = 1$。

\begin{table}[htbp]
\centering
\caption{仿真参数}
\label{tab:param}
\begin{tabular}{c|c}
\hline
参数 & 取值 \\
\hline
基站天线数 $N_t$ & 8, 16, 32 \\
用户数 $K$ & 4, 8, 16 \\
最大发射功率 $P_{\text{max}}$ & 10 dBm \\
信道模型 & 瑞利独立同分布 \\
训练样本数 & 100,000 \\
测试样本数 & 10,000 \\
批次大小 & 256 \\
学习率 & 0.001 \\
\hline
\end{tabular}
\end{table}

\subsection{和速率性能对比}
图2给出了$N_t=16$、$K=8$时不同方法的和速率随信噪比（SNR）变化的曲线。可以看出，DPNet在所有SNR区间均优于ZF-EPA和ZF-WF，在中等SNR（10~20dB）时，DPNet与WMMSE的性能差距小于0.5 bit/s/Hz。高SNR时，DPNet略低于WMMSE（约0.8 dB信噪比损失）。$N_t=32$、$K=16$场景下的趋势相似，这也验证了所提方法适用于大规模MIMO。

\begin{figure}[htbp]
\centering
\includegraphics[width=0.8\linewidth]{example-image}
\caption{不同方法和速率对比（$N_t=16, K=8$）}
\label{fig:sumrate}
\end{figure}

\subsection{计算复杂度与运行时间}
表2对比了各种方法的平均在线计算时间（CPU：Intel i7-12700，PyTorch 2.0）。DPNet的推理时间仅需0.3ms，比WMMSE（30次迭代）快约20倍。ZF-EPA最快，但性能最差。

% 确保表格宽度适配页面
\begin{table}[htbp]
\centering
\caption{不同方法的运行时间对比}
\label{tab:complexity}
\resizebox{\linewidth}{!}{
\begin{tabular}{lcc}
\toprule
方法 & 运行时间/ms & 和速率 (15dB)/bit/s/Hz \\
\midrule
ZF-EPA & 0.05 & 12.3 \\
ZF-WF & 0.12 & 14.1 \\
WMMSE(30 iter) & 6.2 & 18.7 \\
DPNet(本文) & 0.31 & 18.2 \\
\bottomrule
\end{tabular}}
\end{table}

\subsection{鲁棒性分析}
在测试阶段，我们在信道估计中引入高斯误差：$\hat{\mathbf{H}} = \mathbf{H} + \varepsilon \mathbf{E}$，其中$\mathbf{E}$各元素服从标准复高斯分布，$\varepsilon$控制误差强度。图3给出了当$\varepsilon = 0.1$时各方法性能下降情况。WMMSE下降明显（约1.2 dB），而DPNet仅下降0.4 dB，表明神经网络对噪声具有一定的正则化效果。

\begin{figure}[htbp]
\centering
\includegraphics[width=0.8\linewidth]{example-image}
\caption{信道估计误差下的和速率损失（SNR=15dB）}
\label{fig:robustness}
\end{figure}

\section{结论}

本文针对Massive MIMO系统中预编码与功率分配的联合优化问题，提出了一种基于深度学习的端到端解决方案。通过设计复数域神经网络架构和无监督学习框架，实现了从信道状态信息到最优预编码矩阵的直接映射。主要结论如下：

\begin{itemize}
    \item \textbf{性能逼近最优}：所提方法在多种信道条件下均能达到WMMSE算法性能的95%以上，验证了深度学习方法的有效性。
    \item \textbf{计算效率显著提升}：相比传统迭代算法，推理时间降低两个数量级以上，满足实时通信需求。
    \item \textbf{鲁棒性增强}：对信道估计误差具有良好的容忍度，实际应用价值高。
\end{itemize}

未来研究方向包括：1）扩展至多小区协作场景；2）结合能效优化目标；3）考虑硬件实现约束；4）探索联邦学习框架下的分布式优化方案。本研究为智能通信系统的设计提供了新的思路和方法。

\appendix
\section{损失函数梯度推导}
损失函数关于预编码矩阵的梯度可由复数微分法则得到。设$R = \sum_k \log_2(1+\gamma_k)$，则$\frac{\partial R}{\partial \mathbf{w}_k} = \frac{1}{\ln2} \cdot \frac{1}{1+\gamma_k} \cdot \frac{\partial \gamma_k}{\partial \mathbf{w}_k^*}$，其中$\frac{\partial \gamma_k}{\partial \mathbf{w}_k^*} = \frac{\mathbf{h}_k^H \mathbf{h}_k \mathbf{w}_k}{N_k} - \frac{\gamma_k}{N_k} \frac{\partial N_k}{\partial \mathbf{w}_k^*}$，$N_k = \sum_j |\mathbf{h}_k \mathbf{w}_j|^2 + \sigma_k^2$。具体展开略，利用链式法则即可实现自动微分。

\begin{thebibliography}{99}
\bibitem{1} L. Lu, G. Y. Li, A. L. Swindlehurst, et al. An overview of massive MIMO: Benefits and challenges[J]. IEEE J. Sel. Topics Signal Process., 2014, 8(5): 742-758.
\bibitem{2} Q. Shi, M. Razaviyayn, Z. Q. Luo, et al. An iteratively weighted MMSE approach to distributed sum-utility maximization for a MIMO interfering broadcast channel[J]. IEEE Trans. Signal Process., 2011, 59(9): 4331-4340.
\bibitem{3} H. Huang, Y. Yang, J. Yang, et al. Deep learning for precoding in massive MIMO systems[C]//IEEE GLOBECOM Workshops, 2018: 1-6.
\bibitem{4} M. Eisen, A. Ribeiro. Optimal wireless resource allocation with random edge graph neural networks[J]. IEEE Trans. Signal Process., 2020, 68: 2977-2991.
\bibitem{5} C. Trabelsi, O. Bilaniuk, Y. Zhang, et al. Deep complex networks[C]//ICLR, 2018: 1-19.
\end{thebibliography}

\vspace{1em}

\vspace{1em}

\vspace{1em}
{\xiaowu
\noindent \textbf{[作者简介]}\\[-4pt]
\noindent\begin{minipage}[t]{2.5cm}
    \vspace{4pt}
    \rule{2.5cm}{3.5cm}
\end{minipage}\hspace{2pt}
\begin{minipage}[t]{10cm}
    \vspace{4pt}
    秦政（\fontfamily{ptm}\selectfont 1990-），男，重庆人，博士，重庆邮电大学教授、博士生导师，主要研究方向为\fontfamily{ptm}\selectfont MIMO\fontfamily{rm}\selectfont 通信、深度学习与物理层融合。
\end{minipage}
}

\end{document}